%%%%%%%%%%%%%%%%%%%%%
% Documento maestro %
%%%%%%%%%%%%%%%%%%%%%
\documentclass{fime}

%%%%%%%%%%%%%%%%%%%%%%%%%%%%%%%%%%%%%%%%%%%
% Cargando paquetes y definiendo opciones %
%%%%%%%%%%%%%%%%%%%%%%%%%%%%%%%%%%%%%%%%%%%
% Aquí puedes cargar los paquetes que vas a usar. La clase
% fime ya incluye babel, inputenc, graphicx y los de la AMS.
% Cargar un paquete está a tu libertad (y responsabilidad).
\usepackage{hyperref}
    \hypersetup{breaklinks=true,colorlinks=true,
        linkcolor=black,citecolor=black,urlcolor=black}
\usepackage{datetime} % monthname

%%%%%%%%%%%%%%%%%%%%%
% Definiendo campos %
%%%%%%%%%%%%%%%%%%%%%
\def\titulo{Optimal Location of Car Wreck Adjusters}
\def\autor{Luis Alberto Maltos Ortega}
\def\matricula{1390200}
\def\grado{Maestría en Ciencias en Ingenie\'ia de Sistemas}
% En caso de que el grado tenga orientación o especialidad
% llenar el siguiente campo dejando un ESPACIO INICIAL,
% en caso contrario, dejar vacío
\def\orientacion{ }
\def\fecha{\monthname,  \the\year} % Coloca el mes con mayúscula inicial

\def\asesor{Dr. Roger Z. R\'ios}
\def\revisorA{Mar\'ia Ang\'elica Salazar Aguilar}
\def\revisorB{Mar\'ia Guadalupe Villareal Marroqu\'in}
% En el caso de que tu tesis sea de doctorado
% activa la variable cambiándola a \doctoradotrue
% y define tus otros dos revisores
\newif\ifdoctorado\doctoradofalse
\def\revisorC{Nombre del revisor C}
\def\revisorD{Nombre del revisor D}
% El visto bueno siempre va
\def\vobo{Dr. Simón Martínez Martínez}

%%%%%%%%%%%%%%%%%%%%%%%
% Inicia el documento %
%%%%%%%%%%%%%%%%%%%%%%%
\begin{document}

\frontmatter
\pagestyle{main}

%%% Incluye PortillasM si tu tesis es de Maestría
%%% y PortillasD si es de doctorado.
% Portadas (Maestría)

\def\uanl{Universidad Autónoma de Nuevo León}
\def\fime{Facultad de Ingeniería Mecánica y Eléctrica}
\def\depg{Subdirección de Estudios de Posgrado}
\def\snnl{San Nicolás de los Garza, Nuevo León}

%%%%%%%%%%%%%%%%%%%%%%%%
% Primer portada: UANL %
%%%%%%%%%%%%%%%%%%%%%%%%
\thispagestyle{empty}

\begin{scshape}
\begin{center}
	{\Large \uanl} \\[5mm]
	{\large \fime} \\[5mm]
	{\large \depg}
	\vskip 15mm
	\includegraphics[height=55mm]{uanl}
	\vskip 12mm
	\begin{tabular}{p{11cm}}
		\centering
		{\large \titulo}
	\end{tabular}
	\vskip 7mm
	{por}\\[7mm]
	{\large \autor}\\[7mm]
	{como requisito parcial para obtener el grado de}\\[3mm]
	\MakeUppercase{\grado}\\
	\orientacion
	\vfill
	\fecha
\end{center}
\end{scshape}

%%%%%%%%%%%%%%%%%%%%%%%%%
% Segunda portada: FIME %
%%%%%%%%%%%%%%%%%%%%%%%%%
\newpage
\thispagestyle{empty}

\begin{scshape}
\begin{center}
	{\Large \uanl} \\[5mm]
	{\large \fime} \\[5mm]
	{\large \depg}
	\vskip 16mm
	\includegraphics[height=50mm]{fime}
	\vskip 16mm
	\begin{tabular}{p{11cm}}
		\centering
		{\large \titulo}
	\end{tabular}
	\vskip 7mm
	{por}\\[7mm]
	{\large \autor}\\[7mm]
	{como requisito parcial para obtener el grado de}\\[3mm]
	\MakeUppercase{\grado}\\
	\orientacion
	\vfill
	\fecha
\end{center}
\end{scshape}

%%%%%%%%%%%%%%%%%%%%%%%%%%%%%
% Carta del comité de tesis %
%%%%%%%%%%%%%%%%%%%%%%%%%%%%%
\newpage
\thispagestyle{empty}
\enlargethispage{5mm}

\begin{center}
{\bf \large \uanl} \\
{\bf \fime} \\
{\bf \depg}
\end{center}
\vskip 4mm

Los miembros del Comité de Tesis recomendamos que la Tesis <<\titulo>>, realizada por el alumno \autor, con número de matrícula \matricula, sea aceptada para su defensa como requisito parcial para obtener el grado de \grado\orientacion.
\ifdoctorado{\vskip 10mm}\else{\vskip 8mm}\fi

\begin{center}
El Comité de Tesis\\
\ifdoctorado{\vskip 15mm}\else{\vskip 25mm}\fi

\ifdoctorado{%%%
\begin{tabular}{p{37mm}p{21mm}p{12mm}p{21mm}p{37mm}}
	\cline{2-4}
	& \multicolumn{3}{c}{\asesor} & \\
	& \multicolumn{3}{c}{Asesor}  & \\[15mm]
	\cline{1-2} \cline{4-5}
	\multicolumn{2}{c}{\revisorA} & & \multicolumn{2}{c}{\revisorB} \\
	\multicolumn{2}{c}{Revisor}   & & \multicolumn{2}{c}{Revisor}   \\[17mm]
	\cline{1-2} \cline{4-5}
	\multicolumn{2}{c}{\revisorC} & & \multicolumn{2}{c}{\revisorD} \\
	\multicolumn{2}{c}{Revisor}   & & \multicolumn{2}{c}{Revisor}   \\[2mm]
	& \multicolumn{3}{c}{Vo. Bo.} & \\[14mm]
	\cline{2-4}
	& \multicolumn{3}{c}{\vobo}   & \\
	& \multicolumn{3}{c}{\depg}   & \\ &&&&
\end{tabular}
}\else{%%%
\begin{tabular}{p{37mm}p{21mm}p{12mm}p{21mm}p{37mm}}
	\cline{2-4}
	& \multicolumn{3}{c}{\asesor} & \\
	& \multicolumn{3}{c}{Director}  & \\[19mm]
	\cline{1-2} \cline{4-5}
	\multicolumn{2}{c}{\revisorA} & & \multicolumn{2}{c}{\revisorB} \\
	\multicolumn{2}{c}{Revisor}   & & \multicolumn{2}{c}{Revisor}   \\[2mm]
	& \multicolumn{3}{c}{Vo. Bo.} & \\[17mm]
	\cline{2-4}
	& \multicolumn{3}{c}{\vobo}   & \\
	& \multicolumn{3}{c}{\depg}   & \\ &&&&
\end{tabular}
}\fi%%%

\vfill

\snnl, \MakeLowercase{\fecha}

\end{center}

% Dedicatoria

\thispagestyle{empty}
\vspace*{17mm}

\begin{flushright}
\begin{itshape}

Aquí puedes poner tu dedicatoria\\
si es que tienes una.\bigskip\bigskip

Si no tienes una, puedes borrar\\
la línea \verb+% Dedicatoria

\thispagestyle{empty}
\vspace*{17mm}

\begin{flushright}
\begin{itshape}

Aquí puedes poner tu dedicatoria\\
si es que tienes una.\bigskip\bigskip

Si no tienes una, puedes borrar\\
la línea \verb+% Dedicatoria

\thispagestyle{empty}
\vspace*{17mm}

\begin{flushright}
\begin{itshape}

Aquí puedes poner tu dedicatoria\\
si es que tienes una.\bigskip\bigskip

Si no tienes una, puedes borrar\\
la línea \verb+\include{Dedicatoria}+ en el\\
archivo \texttt{MiTesis.tex} pues no es obligatoria.

\end{itshape}
\end{flushright}

+ en el\\
archivo \texttt{MiTesis.tex} pues no es obligatoria.

\end{itshape}
\end{flushright}

+ en el\\
archivo \texttt{MiTesis.tex} pues no es obligatoria.

\end{itshape}
\end{flushright}



\tableofcontents
\listoffigures
\listoftables

%Acknowledgements

\chapter{Acknowledgments}
\markboth{Acknowledgments}{}

%Aquí puedes poner tus agradecimientos.
%(No olvides agradecer a tu comité de tesis,
I want to especially thank
Prof. Roger Z. Rios
For giving me the opportunity to work with him
and be my advisor,
for guiding me in my academic development
and for the effort and support he gave me
I also would like to thanks him
for his patience,
without his support
this work would not have been possible. 

Thanks to my thesis committee
Prof. Ma. Angelica Salazar Aguilar
and Prof. Ma. Guadalupe Villareal Marroquin,
for their comments and suggestions
that made throughou his work,
and for help me
to complete this work,

%a tus profesores,
to all professors of the PISIS
for their courses so complete and enriching,

%a la facultad
to the FIME
for the facilities granted
and ease the paperwork,

%y a CONACyT
%en caso de que hallas sido beneficiado con una beca).
to Consejo Nacional de Cíencia y Tecnología (CONACyT)
for the Graduate Fellowship
and the research project grant 2011-1-166397
with which it was possible the realization of my studies
and attendance at national conferences
to present part of this work.

%Resumen

\chapter{Abstract}
\markboth{Abstract}{}

{\setlength{\leftskip}{10mm}
\setlength{\parindent}{-10mm}

\autor.

Candidato para obtener el grado de \grado\orientacion.

\uanl.

\fime.

Título del estudio: \textsc{\titulo}.

\noindent Número de páginas: \pageref*{lastpage}.}

%%% Comienza a llenar aquí
\paragraph{Problem Description:}
When a traffic accident occurs
in cities with a large traffic flow
the roads surrounding the crash site
are affected by traffic congestions.
In some countries, such as Mexico,
even small accidents are troublesome
due to the fact
that a claim adjuster
from the car insurance company
must arrive to the site
and document the accident
before the vehicle may be moved
as required by law.
Thus,
in this particular setting,
the adjusters location
becomes a key factor in providing timely service.

\paragraph{Objectives and method of study:}
%Aquí debes poner tus objetivos y métodos de estudio.
The aim of this thesis is to
\begin{itemize}
\item Provide quantitative tools
with scientific support
for the optimal location
of car wreck adjusters.
\item Develop
adequate mathematical models
for representing
some of the important company concerns.
Design and develop efficient solution techniques
for handling real-world instances of the problem.
\item Assess
the performance of the proposed techniques
based on an appropriate experimental design.
\end{itemize}

\paragraph{Contributions and Conclusions:}
%Y aquí tus contribuciones y conclusiones.
%(También es parte del formato).
The most important contributions of this work
are the following:
\begin{itemize}
\item Two integer programming models
  are introduced
  with the aim of
  minimizing average response time.
\item The second model
  takes into account
  special features
  to reduce the number of variables.
\item A scatter search heuristic
  was designed, implemented, and tested
  on a wide set of instances
  with very good results.
\end{itemize}

\bigskip\noindent\begin{tabular}{lc}
\vspace*{-2mm}\hspace*{-2mm}Firma del asesor: & \\
\cline{2-2} & \hspace*{1em}\asesor\hspace*{1em}
\end{tabular}




\mainmatter
\pagestyle{fime}

%%% Haz un documento para cada capítulo
\chapter{Introduction}
The main idea behind this thesis
is to develop models and methods
for improving the service offered by car insurance agents,
helping them arrive to accident sites sooner,
and determining the number of adjusters required
to perform the service within the desired standards.

\section{Problem Statement}
A car insurance company
in countries such as Mexico
faces the daily issue
of how to locate
their insurance agents (adjusters),
in such a way
that they provide
the best possible service
to their customers.
Typically
the quality of service
is measured
by how quickly
the adjusters arrive
to the places
where an incident has occurred.

This problem
can be seen
as an emergency location system, % add cite
in which
it is desired to ``cover''
the territory under study
(city intersections)
by a fixed number of adjusters
so as to reduce
the average time
of arrival to accidents.
Another issue
that makes the problem
even more complex
is the fact
that accidents occur randomly
and decisions
on where to place the adjusters
must be made before
accidents occur.

\section{Background}
This is a practical problem that has not been studied,
ya que este sistema de emergencia,
no se presenta en muchos paises,
sin embargo, hay sistemas similares,
tal es el caso de ambulancias, bomberos, patrullas,
para los cuales existen una gran cantidad de articulos al respecto,
pero las caracteristicas tratadas en ellos,
difieren de nuestro problema practico,
por los siguientes factores
tipo de espacio,
radio cobertura,
multiperiodo,
tiempo de servicio,
tipo de llegadas

Berman propone un problema de localizacion para sistemas de emergencia,
y lo llama The Queue Stochastic p-Median Problem.

\section{Motivation}
When a car accident occurs,
traffic congestion starts to pile up.
This is because
customers are not allowed
to move their vehicles until
the adjuster arrives.
The adjuster must record and determine
the causes of the accident,
in order to
move the car from the accident area
and restore the flow.

\section{Objectives}
The aim of this work
is to support decision-making,
regardin location of insure agents
to attend car wrecks.

The main goal is
to determine the optimal bases (locations)
for placing the insurance company adjusters,
so as to minimize
the average or maximum response time
from customer calls
when accidents occur.

\section{Organization}
This thesis is organized as follows:
Chapter \ref{ch:RelatedWork}
presents an brief literature review
of diverse approaches of Emergency Service Systems (ESS),
starting with the simple deterministic models
and ending with the hypercube model and simulation models.
The framework is discussed in Chapter \ref{ch:Framework}
the problem statement and two models are presented.
Chapter \ref{ch:Heuristic} describes the heuristic method proposed
and their components.
Chapter \ref{ch:Experiments} contains
the experiments made to the models,
and to the heuristics to tune them.
Conclusions, contributions, 
and directions for future research
are highlighted in Chapter \ref{ch:Conclusions}.
Additionally, Appendix \ref{ch:Jarvis}
contains a summary of how to approximate the hypercube results,
and together with Appendix \ref{ch:Berman}
completes the explanation of how to estimate
the response times of a solution.
%cedure discussed in Chapter 5
%Chapter 6 describes
%Chapter 7 contains a
%Conclusions, contributions, 
%and directions for future research are highlighted in Chapter 8
%Additionally, Appendix A


\chapter{Models}

% Validating and applying a model 
% for locating emergency medical vehicles in Tucson, AZ
% Goldberg 1990

\section{Model A}

This model is based
on the model
proposed by Goldberg \cite{goldberg1990validating},
and it contains assignment variables
for all possible orders.

The following assumptions are considered in the model:
\begin{itemize}
\item The probability that an adjuster is busy
  is $\rho$ and is unaffected by the state of the system
\item There is a strict ordering of the bases preferred for each zone
  that does not depend
  on the current state of the system
\item All calls are answered
  by an adjuster originating from its base,
  not en route back to the base
\item The arrival of calls to the system
  follows a stationary distribution
\item The model is presented
  using a 0-queue assumption.
\end{itemize}

Sets and indexes:
\begin{itemize}
\item \textit{n} number of demand points
\item \textit{m} number of potential sites to locate adjusters/facilities
\item \textit{p} number of available adjusters
\item \textit{i} index for demand points;
  $i \in V = \{1,2,\ldots,n\}$ 
\item \textit{j} index for potential site for adjusters/facilities
  $j \in W = \{1,2,\ldots,m\}$
\item $k$ possible order;
  $k \in K = \{1,2\ldots,p\}$
\item $S_{ij} = \{
  r \in W | \mbox{ site } r \mbox{ is preferred before site } j
  \mbox{ for demand point } i
  \}$
\end{itemize}

Parameters:
\begin{itemize}
\item $\rho$ is the utilization of each adjuster
\item $t_{ij}$ is the expected travel time
  between demand point \textit{i}
  and potential site \textit{j}.
\item $h_{ij}^{k}$ is the probability
  that $j$ serves point $i$
  given that
  is the $k$-th preferred.
\end{itemize}

Variables:
\begin{itemize}
\item $x_j =
  \begin{cases} 
    1 & \mbox{if an adjuster is at potential site } j \\
    0 & \mbox{otherwise}
  \end{cases}$
\item $y_{ij}^{k} =
  \begin{cases} 
    1 & \mbox{if adjuster in } j \mbox{, is the }
    k\mbox{-th to cover demand point } i \\
    0 & \mbox{otherwise}
  \end{cases}$
\end{itemize}

Model
{\small
  \begin{equation}
    \min \, \sum_{j=1}^{m}{
      \sum_{k=1}^{p}{
        \sum_{i=1}^{n}{
          h_{ij}^{k}t_{ij}y_{ij}^{k}
        }
      }
    }
  \end{equation}
}
Minimize the average expected response time
\begin{align}
  \label{eq:2}
  \sum_{j \in W}{x_j}
  & = p
\end{align}
Only locate $p$ adjusters
\begin{align}
  \label{eq:3}
  \sum_{j \in W}{y_{ij}^{k}}
  & = 1
  & i \in V, k
  &\in K
\end{align}
Each demand point $i$ is covered by an adjuster on each order $k$
\begin{align}
  \label{eq:4}
  y_{ij}^{k}
  & \leq x_j
  & i \in V,j \in W, k
  &\in K
\end{align}
Assigned only if there is an adjuster located at $j$
\begin{align}
  \label{eq:5}
  \sum_{k = 1}^{p}{
    y_{ij}^{k}
  }
  & \leq x_j
  & i \in V, j 
  & \in W
\end{align}
Only assign in one order for each adjuster located 
\begin{align}
  y_{ij}^{k} 
  & \leq \sum_{r\in S_{ij}}{y_{ir}^{k-1}}
  & i \in V,j \in W, k
  & \in K/\{1\}
\end{align}
Assign \textit{j} to cover \textit{i} in order $k$
only if
the assign of order $k-1$
was made for some $r \in S_{ij}$
\begin{align}
  x_{j}
  & \in \{0,1\}
  & j 
  & \in W \nonumber
  \\
  y_{ij}^{k}
  & \in \{0,1\}
  &  i \in V,j \in W,k
  &\in K \nonumber
\end{align}

%This model was tested 
%with a benchmark instance of 30 points, 
%which sets $V$ and $W$ are the same.
%For values of $p = 1,\ldots,7$
%the optimal solution was found in less than one minute.
%
%For other benchmark instances
%was impossible to find the solution
%because they have more than 300 points,
%and the proposed model grows to over one million variables.
%  For some large instances ($n = 50$ , $m = 30$)%
\todo[inline]{Add Scope}
\todo[inline]{Computational Results}

% A model for the Stochastic Queue p-Median Problem
% Created by Luis Maltos & Roger Rios
% 2015

\section{Model B}

This model
was developed with the idea
that is unlikely that the farthest adjusters
serve demand points
on cases where the system does not become congested.
In these cases
we can make the assumption that
the probability of being served by the $\ell$-th 
adjuster is almost zero,
where $\ell$ is large enough but less than $p$.

Parameters:
\begin{itemize}
\item $M$ is a large integer
\item $\ell$ the number of allowed adjusters per demand point
\item $a_{ik}$ the $k$-th preferred location server
  regarding the point $i$.
\end{itemize}

Variables:
\begin{itemize}
\item $z_j$ the number of adjusters placed at site $j$
\item $y_{ij}^k$ if adjuster in $j$,
  is the $k$-th to cover demand point $i$
\end{itemize}

The objective,
and the constraints \ref{eq:2}-\ref{eq:5} 
are practically the same as in model A,
with the difference that
binary variables $x_j$ from Model A
are replaced by integer variables $z_j$
inspired by the results of Berman \cite{berman1987stochastic},
and the addition of the following binary variables
\begin{itemize}
\item $u_{ij} = 
  \begin{cases}
    1 & \mbox{if the number of adjusters between } i 
    \mbox{ and } j \mbox{, inclusive, is less than } \ell \\
    0 & \mbox{otherwise}
  \end{cases}$
\item $v_{ij} = 
  \begin{cases} 
    1 & \mbox{if the number of adjusters between } i
    \mbox{ and } j \mbox{, is less than } \ell - 1 \\
    0 & \mbox{otherwise}
  \end{cases}$
\end{itemize}

\begin{align}
  \sum_{r = 1}^{k \mid a_{ik}=j}{
    z_{a_{ir}}
  }
  + (p-\ell) u_{ij}
  & \leq p
  & i \in V,
  & j \in W
\end{align}
Checks if
the number of adjusters between $i$ and $j$,
inclusive,
is less or equal than $\ell$
\begin{align}
  \sum_{r = 1}^{k \mid a_{ik}=j}{
    z_{a_{ir}}
  }
  + M u_{ij}
  & \geq \ell+1
  & i \in V,
  & j \in W
\end{align}
Checks if
the number of adjusters between $i$ and $j$,
inclusive,
is greater than $\ell$
\begin{align}
  \sum_{k = 1}^{\ell}{
    y_{ij}^{k}
  }
  + M (1 - u_{ij})
  & \geq z_j
  & i \in V,
  & j \in W
\end{align}
Assign $z_j$ times $j$ to $i$ if $u_{ij} = 1$
\begin{align}
  \sum_{r = 1}^{k \mid a_{i(k+1)}=j}{
    z_{a_{ir}}
  }
  + (p-(\ell-1)) v_{ij}
  & \leq p
  & i \in V,
  & j \in W
\end{align}
Checks if
the number of adjusters between $i$ and $j$,
is less than $\ell$
\begin{align}
  \sum_{r = 1}^{k \mid a_{i(k+1)}=j}{
    z_{a_{ir}}
  }
  + M v_{ij}
  & \geq \ell
  & i \in V,
  & j \in W
\end{align}
Checks if
the number of adjusters between $i$ and $j$,
is greater or equal than $\ell$
\begin{align}
  \sum_{k=1}^{k}{
    y_{ij}^{k}
  }
  + M (1 - v_{ij} + u_{ij})
  & \geq \ell
  - \sum_{r = 1}^{k \mid a_{i(k+1)}=j}{
    z_{a_{ir}}
  } 
  &  i \in V, j 
  & \in W
\end{align}
Assign $j$ to $i$
the times remaining to complete $\ell$ assignments
\begin{align}
  \sum_{k=1}^{k}{
    y_{ij}^{k}
  }
  - M (1 - v_{ij} + u_{ij})
  & \leq \ell 
  - \sum_{r = 1}^{k \mid a_{i(k+1)}=j}{
    z_{a_{ir}}
  }
  & i \in V,
  & j \in W
\end{align}
Assign $j$ to $i$
the times remaining to complete $\ell$ assignments
\begin{align}
  y_{ij}^{k}
  & \leq u_{ij}
  + v_{ij}
  & i \in V,
  & j \in W
\end{align}
Assign $j$ to $i$
only if
is in the first $\ell$ adjusters near $i$
\begin{align}
  z_j
  & \in \{0,1,\ldots,p\}
  & j
  & \in V \nonumber
  \\
  y_{ij}^{k} 
  & \in \{0,1\}
  & i\in V,j\in W,k
  & \in I \nonumber
  \\
  u_{ij},v_{ij}
  & \in \{0,1\}
  & i \in V,j
  & \in W \nonumber
\end{align}



\appendix
%%% Haz un documento para cada apéndice
\chapter{Este es un apéndice}

\section{Citas bibliográficas}

En princicpio tienes total libertad de incluir tu bibliografía con el entorno {\tt thebibliography} nativo de \LaTeX{} o mediante la herramienta \textsc{Bib}\TeX. En caso de que optes por esto último (recomendado), puedes usar alguno de los archivos {\tt mighelbib.bst} o {\tt mighelnat.bst} incluidos en el paquete {\tt Tesis-FIME}, pues sus diseños están basados en el estilo bibliográfico estándar del español, además de que armoniza con el estilo de tesis provisto por {\tt fime.cls}.

El estilo bibliográfico {\tt mighelbib} es numérico, es decir cita con un número entre corchetes, por ejemplo una cita \verb+\cite{Dan82}+ genera una etiqueta del tipo [13], mientras que el estilo {\tt mighelnat} es tipo autor-año y requiere que el paquete {\tt natbib} sea cargado (sin opciones) para su correcto funcionamiento, cita con el apellido del autor y el año, por ejemplo una cita \verb+\citet{Dan82}+ genera una etiqueta del tipo Dantzig (1982), mientras que una cita \verb+\citep{Dan82}+ genera una etiqueta del tipo (Dantzig, 1982).

Como muestra del estilo, unas citas: un libro clásico de programación lineal \cite{Baz04} y un documento histórico \cite{Dan82}. Para saber un poco más del uso de \textsc{Bib}\TeX, se puede leer \cite{Mat11}.

\section{Comillas}

El objetivo de esta sección era provocar otra página para que se vea el encabezado. Pero aprovechamos para decir que la clase {\tt fime.cls} carga el paquete {\tt babel} con la opción {\tt spanish}, por lo que cambiará automáticamente los dobles signos $<<$ y $>>$ por << y >>. Estas comillas angulares son las correctas en el idioma español, y son las que se usan en la clase {\tt fime.cls}, por lo que se sugiere sean las usadas en el texto cada que quieras <<entrecomillar>> algo.



\backmatter
\pagestyle{main}

%%% Aquí va la bibliografía, puedes usar el entorno de LaTeX (thebibliography)
%%% o la herramienta BibTeX. En caso de que optes por BibTeX, puedes usar
%%% alguno de los archivos de estilo (mighelbib.bst o mighelnat.bst) incluidos
%%% en el paquete, cuyos diseños armonizan con el diseño de tesis provisto por
%%% fime.cls. Para muestra, basta un botón:
\bibliographystyle{mighelbib}
\bibliography{MiBiblio}

\label{lastpage}
%Autobiografia

\chapter*{Resumen autobiográfico}
\thispagestyle{empty}

\begin{center}
\autor

Candidato para obtener el grado de\\
\grado\\
\orientacion\bigskip

\uanl\\
\fime\bigskip

Tesis:\\
\textsc{\large\titulo}
\end{center}\bigskip

%Aquí va tu historia
Aquí va tu historia. Recuerda que debe incluir: lugar y fecha de nacimiento, nombre de los padres, escuelas y universidades en las que se graduó después de la preparatoria, títulos o grados obtenidos (no incluir los estudios que se están concluyendo), experiencia profesional y organizaciones profesionales a las que pertenece (no incluir lista de publicaciones).


\end{document}
